\section{Preliminares}\label{sec:preliminares}

\subsection{Notación}

Indexamos filas y columnas por $\{0,\dots,n-1\}$. La casilla $(r,c)$ está en la fila $r$, la
columna $c$, la diagonal $r-c\in\{-(n-1),\dots,n-1\}$ y la antidiagonal
$r+c\in\{0,\dots,2n-2\}$. Llamamos \emph{líneas} a filas, columnas, diagonales y antidiagonales,
y $R,C,D,A$ a sus cuatro familias, con $n$, $n$, $2n-1$ y $2n-1$ elementos. El tablero es así un
hipergrafo $4$-partito y $4$-uniforme cuyos vértices son las líneas y cuyas aristas son las
casillas, cada una con sus cuatro líneas. Dos líneas distintas están juntas en a lo sumo una
casilla (\emph{codegrado uno}). Una colocación legal es un emparejamiento de este hipergrafo.

Para $P\in\Om_h(n)$, $R_\vac(P)$ y $C_\vac(P)$ son las $h$ filas y las $h$ columnas vacías, y
$X_h(P)$ es el número de cruces vacíos disponibles. Para coordenadas continuas usamos los centros
$\bar r=(r+\frac12)/n$, $\bar c=(c+\frac12)/n$; una diagonal o antidiagonal recibe la coordenada
que inducen los centros de sus casillas, $\bar r-\bar c\in[-1,1]$ o $\bar r+\bar c\in[0,2]$.

Salvo indicación, $n\to\infty$ con $h$ (o $K$) fijo; $c,C$ denotan constantes positivas que
pueden cambiar de una aparición a otra y cuya dependencia de parámetros fijos se indica con
subíndices. Para leyes $\mu,\nu$ en un conjunto finito,
$\KL{\mu}{\nu}=\sum\mu\log(\mu/\nu)$ es la entropía relativa (infinita si
$\mu\not\ll\nu$). Usaremos que es no negativa y que no aumenta al aplicar una misma función a
ambas leyes (procesamiento de datos)~\cite{CoverThomas}.

\subsection{Interpretación de \texorpdfstring{$f_h$}{fh}}

En la Figura~\ref{fig:cruces-vacios} contamos los cruces en una colocación
particular. Al sumar el mismo contador sobre toda la capa obtenemos la
identidad siguiente.

\begin{lema}\label{lem:interpretacion}
Para $1\le h\le n$, $\E_{\mu_h}[X_h]=h^2f_h(n)$. En particular, $f_1(n)$ es la probabilidad de
que una colocación uniforme de $\Om_1(n)$ pueda completarse a una solución.
\end{lema}

\begin{proof}
Contamos los pares $(P,P')$ con $P\in\Om_h$, $P'\in\Om_{h-1}$ y $P\subset P'$. Cada $P'$ tiene
$n-h+1$ reinas y al retirar cualquiera de ellas se obtiene un $P$ distinto: hay
$(n-h+1)N_{h-1}$ pares. Recíprocamente, añadir una reina a $P$ da una colocación legal si y
sólo si la reina ocupa un cruce vacío disponible: hay $X_h(P)$ extensiones. Así
$N_h\E_{\mu_h}X_h=(n-h+1)N_{h-1}$, que es la identidad buscada. Si $h=1$, $X_1\in\{0,1\}$ y
$X_1(P)=1$ exactamente cuando $P$ se completa.
\end{proof}

\subsection{El problema variacional y el perfil}\label{subsec:variacional}

La doble cuenta anterior es exacta para cada tamaño. Para estudiar su límite
necesitamos una descripción del tablero que conserve las restricciones de fila,
columna y diagonales al pasar a coordenadas continuas.

Para una masa $\rho\in(0,1]$, una densidad \emph{admisible} es una función medible
$\nu\ge0$ en $[0,1]^2$ con $\int\nu=\rho$ cuyas proyecciones (o marginales) sobre las cuatro familias,
\[
 m_R(r)=\int_0^1\nu(r,c)\,dc,\quad m_C(c)=\int_0^1\nu(r,c)\,dr,\quad
 m_D(d)=\int_{r-c=d}\nu,\quad m_A(s)=\int_{r+c=s}\nu,
\]
no exceden uno (las dos últimas integran a lo largo de la fibra, parametrizada por $r$). Sus
vacancias son $u_f=1-m_f$ en $I_R=I_C=[0,1]$, $I_D=[-1,1]$, $I_A=[0,2]$; así
$\int u_R=\int u_C=1-\rho$ y $\int u_D=\int u_A=2-\rho$. El funcional de entropía es
\[
 s[\nu]=-3\rho-\int\nu\log\nu-\sum_{f\in\{R,C,D,A\}}\int_{I_f}u_f\log u_f ,
\]
y $s^*(\rho)$ es su supremo sobre las densidades admisibles de masa $\rho$; ponemos
$s_1=s^*(1)$. Para $0<\rho<1$ el óptimo $\nu$ es positivo y satisface la ecuación de Euler
$\nu=\lambda_\rho\,u_Ru_Cu_Du_A$, con una constante $\lambda_\rho>0$, la \emph{actividad}~\cite[Teorema~3.1 y Lema~3.3]{Perfil}.

Con $x=1-\rho$ se reescala: $r_R=u_R/x$, $r_C=u_C/x$ y $b_x=\lambda_\rho x^2$, la actividad reescalada. Entonces
\begin{equation}\label{eq:perfil}
 \nu=b_x\,r_Rr_Cu_Du_A,\qquad \int r_R=\int r_C=1,
\end{equation}
y las ecuaciones de Euler se escriben
\begin{equation}\label{eq:euler}
 r_R=\frac1{x+b_xJ_R},\quad r_C=\frac1{x+b_xJ_C},\quad
 u_D=\frac1{1+b_xJ_D},\quad u_A=\frac1{1+b_xJ_A},
\end{equation}
donde cada $J_f$ integra a lo largo de la fibra de $f$ el producto de los otros tres factores;
por ejemplo $J_R(r)=\int_0^1r_C(c)u_D(r-c)u_A(r+c)\,dc$. En efecto, la masa de la fila $r$ es
$1-xr_R(r)=b_xr_R(r)J_R(r)$. Integrando~\eqref{eq:perfil} se obtiene la expresión~\eqref{eq:F}
de $F(x)=(1-x)/b_x$. Escribimos $\nu_x$ y $(b_x,r_{R,x},r_{C,x},u_{D,x},u_{A,x})$ cuando haga
falta indicar $x$; el subíndice es entonces la fracción de huecos, no la masa
(en~\cite{Perfil} las densidades se indexan por la masa: nuestra $\nu_x$ es allí $\nu_{1-x}$). La única
excepción es la densidad de masa uno, que corresponde a $x=0$ y se denota $\nu_1$, como
en~\cite{Perfil}.

\begin{teorema}[El perfil cerca del tablero lleno~\cite{Perfil}]\label{thm:perfil}
Existen constantes $0<\kappa\le1$, $0<\underline r\le M<\infty$, $0<b_-\le b_+<\infty$ y
$L_*<\infty$ tales que, para todo $x\in(0,1/2]$, el perfil óptimo de masa $1-x$~\cite[Teorema~3.1]{Perfil}, reescalado como
en~\eqref{eq:perfil}, satisface~\eqref{eq:euler} y:
\begin{enumerate}[(a)]
\item $b_-\le b_x\le b_+$, $\ \underline r\le r_{R,x},r_{C,x}\le M$ y
$\kappa\le u_{D,x},u_{A,x}\le1$;
\item $r_{R,x},r_{C,x},u_{D,x},u_{A,x}$ son Lipschitz en sus intervalos con constante a lo sumo
$L_*$, y $\nu_x$ es acotada y Lipschitz, uniformemente en $x$;
\item cuando $x\downarrow0$ convergen uniformemente a límites $(b_0,r_R,r_C,u_D,u_A)$ que
satisfacen~\eqref{eq:euler} con $x=0$ y las mismas cotas.
\end{enumerate}
En particular, $\nu_1=b_0r_Rr_Cu_Du_A$ es una densidad continua en $[0,1]^2$, de masa uno, con
marginales de filas y columnas iguales a uno y vacancias diagonales $u_D=1-m_D$ y $u_A=1-m_A$, de
modo que $\int u_D=\int u_A=1$; $\nu_1$ es el único maximizador de $s$ entre las densidades
admisibles de masa uno, así que $s[\nu_1]=s^*(1)=s_1$; y $F$ es continua en $0$ con
$F(0)=1/b_0>0$.
\end{teorema}

Las afirmaciones (a) y (c) son las Proposiciones~4.2, 5.2 y~5.3 de~\cite{Perfil}; que el
límite $\nu_1$ sea el único maximizador de masa uno es su Proposición~5.1 junto con la~5.2; la
regularidad Lipschitz uniforme (b) es su Teorema~7.7, y la continuidad de $F$ en $0$, con
$F(0)=1/b_0$, es su Teorema~8.1. El Apéndice~\ref{app:perfil} resume esas pruebas y da valores
explícitos de las constantes. No se usa ninguna identificación del perfil con leyes discretas de colocaciones:
el Teorema~\ref{thm:perfil} es un enunciado sobre el problema continuo.
